\documentclass[10pt,a4paper]{amsart}
\usepackage[margin=19mm]{geometry}
\usepackage{amsmath,amssymb,mathtools}
\usepackage[scr=rsfs,cal=euler]{mathalfa}
\usepackage{xfrac}
\usepackage[unicode,hidelinks,bookmarks=false]{hyperref}
\newcommand{\ie}{i.e.\ }
\newcommand{\cf}{cf.\ }
\newcommand{\resp}{respectively\ }
\newcommand{\Subsection}{\subsection}
\providecommand{\nobreakdash}{\nobreak\hbox{-}}
\setlength{\emergencystretch}{2em}
\multlinegap0pt
\usepackage{tikz}
\usepackage{xcolor}
\usepackage{amssymb}
\newtheorem{lemma}{Lemma}
\newcommand{\Z}{\mathbb Z}
\def\le{\leqslant}
\title{Convergence of sparse square-summable NLFT}
\author{Sergey A. Denisov}
\date{Source: arXiv:2606.06661v2 (CC BY 4.0)}
\begin{document}
\maketitle

\noindent\emph{Source and licence.} Convergence of sparse square-summable NLFT. Version arXiv:2606.06661v2 (CC BY 4.0), distributed by arXiv under the Creative Commons Attribution 4.0 International licence, http://creativecommons.org/licenses/by/4.0/, as recorded in the arXiv licence metadata for this exact version and shown on the abstract page https://arxiv.org/abs/2606.06661v2. Full licence notice: \texttt{licenses/arxiv-2606.06661-CC-BY-4.0.txt}.\par\smallskip

\noindent\textit{Editorial scope.} Counted unit: the article's lemma suv (source0.tex lines 278--279). The article's complete original proof of it is delivered with it (source0.tex lines 280--281, one \texttt{proof} environment of 2 lines). No mathematics is written, repaired or re-derived: the statement and the proof are the article's own lines, in the article's own notation, and no retained line is substituted or re-flowed.  Retained uncounted as the source-local context the proof needs: lines 246--248; lines 270--272.  Nothing was omitted from any retained span, so the package carries no omission record.  No retained line contains a citation, so the wrapper carries no bibliography and no bibliography entry.  No retained line contains a figure environment, an image-inclusion command or drawing code, so no image is delivered and the package declares no asset dependency.  Editorial formatting only, in addition: an AMS \texttt{amsart} preamble that restates the article's own declarations and macros used by the retained lines, namely the environments \texttt{lemma} on separate counters, together with the standard packages those lines need.  The retained environments print in this document as Lemma 1 in order of appearance, whereas the article prints the same items with the numbering of its own full document.  The complete native source of the article as distributed in the arXiv e-print for this exact version is bundled unchanged as \texttt{sources/202200/source0.tex}.
\begin{equation}\label{sd1}
\alpha_{j}=\alpha_{j-1}+{\delta}_j z^{n_j}\overline{\beta}_{j-1}, \quad \beta_{j}=\beta_{j-1}+\delta_j z^{n_j}\overline{\alpha}_{j-1}\,.
\end{equation}

\begin{equation}\label{le}
\Lambda_j\subset \Lambda_{j-1}\cup (n_j-\Sigma_{j-1}), \quad \Sigma_j\subset \Sigma_{j-1}\cup (n_j-\Lambda_{j-1})\,
\end{equation}

\begin{lemma}\label{suv} If $\alpha_j=:\sum_{0\le s}x_s^{(j)}z^s, \beta_j=:\sum_{0\le s}y_s^{(j)}z^s$, then for fixed $s$ we get $x_s^{(j)}=x_s$ and $y_s^{(j)}=y_s$ for large enough $j$. Moreover, $x:=\{x_s\}\in \ell^2(\Z^+)$ and $y:=\{y_s\}\in \ell^2(\Z^+)$.
\end{lemma}

\begin{proof}The first claim follows from   \eqref{sd1} and \eqref{le} which show that $\{x_s\}$ and $\{y_s\}$ are defined recursively. The second  claim follows  from the previous lemma.
\end{proof}

\end{document}
